
\subsubsection{Influence Functions and Data Attribution}

Influence functions \citep{koh2017understanding} estimate training example influence by approximating leave-one-out retraining through the inverse Hessian-vector product (IHVP). While theoretically principled, IHVP computation is prohibitively expensive for modern deep networks. Subsequent work developed three main approximation strategies:

\textbf{First-order methods} bypass curvature entirely. TracIn \citep{pruthi2020estimating} approximates influence via gradient dot-products at training checkpoints, achieving O(1) complexity per example pair. While computationally attractive, TracIn ignores second-order information.

\textbf{Curvature approximation methods} factorize or approximate the Hessian. EK-FAC \citep{grosse2023studying} extends KFAC's Kronecker factorization to influence estimation, scaling to 52B parameter LLMs with 0.85 Spearman correlation. However, their evaluation focused exclusively on decoder-only architectures (LLaMA-2), leaving encoder behavior unexplored.

\textbf{Random projection methods} project gradients to lower-dimensional spaces. TRAK \citep{park2023trak} uses random feature regression to estimate datamodel coefficients. Their evaluation included BERT and CLIP, but tested each architecture in isolation without controlled comparison.

A critical gap exists: no prior work has systematically compared these methods across architectures at matched conditions. Each method was validated on its ''home'' architecture---EK-FAC on decoders, TracIn and TRAK on various models---without controlling for architecture-specific effects.

\subsubsection{Transformer Architectures and Attention Structure}

The transformer architecture \citep{vaswani2017attention} uses attention mechanisms that fundamentally differ between encoder and decoder variants. Encoder-only models \citep{devlin2019bert} employ bidirectional attention where each token attends to all positions. Decoder-only models \citep{radford2019language} use causal attention masks that prevent attending to future tokens.

This architectural difference has implications for gradient computation. Bidirectional attention creates O($n^{2})$ dense attention weights, while causal attention computes only O($n^{2}/2)$ non-zero weights. Prior work has noted that causal structure creates ''block-diagonal-ish'' attention Jacobians \citep{grosse2023studying}, but the downstream effect on attribution accuracy remained unexplored.

\subsubsection{Contribution}

This work provides the first systematic matched cross-architecture comparison of data attribution methods. By controlling for model size (both ~110-125M parameters), depth (both 12 layers), task (SST-2 classification), and training procedure (identical hyperparameters), the effect of attention structure on attribution performance is isolated.
